
\subsubsection{2.1 Weight-Space Property Prediction}

Unterthiner et al. (2020) demonstrated that properties of trained neural networks are predictable from their weights with high fidelity on a zoo of over 120,000 models. Their flat-MLP baseline concatenates all weight matrices into a single vector and applies a standard MLP regression head. Eilertsen et al. (2020) extended this direction with meta-classifiers operating on weight snapshots (320K snapshots across 16K networks). Both works operate in the in-distribution setting with implicitly consistent neuron ordering and do not evaluate behavior under permutation stress.

Schurholt et al. (2021) introduced hyper-representations using contrastive learning with permutation augmentation, demonstrating that augmenting with random neuron permutations during training can improve generalization of weight-space encoders. More recent work has extended weight-space property prediction to graph-structured model representations \citep{WSKAN2026} and large-scale benchmarking across diverse model zoos \citep{Schurholt2025}.

\subsubsection{2.2 Permutation-Equivariant Architectures for Weight Spaces}

Zhou et al. (2023) introduced Neural Functional Transformers (NFT), which apply permutation-equivariant attention layers to neural network weight representations. NFT represents each neuron as a token and applies multi-head attention where equivariance is enforced across the neuron dimension. The equivariance theorem (Theorem 1 of Zhou et al., 2023) guarantees that permuting neurons permutes the attention outputs correspondingly. NFT was evaluated on INR classification tasks but not on model zoo property prediction.

Navon et al. (2023) developed Deep Weight Space Networks (DWSNets), which achieve equivariance for CNN weight spaces. DWSNets encounters runtime shape mismatches with FC-MLP weight vectors and was not included as a baseline.

Subsequent work has extended equivariant weight-space representations to graph metanetworks \citep{Kofinas2024} and architecture-agnostic encoders \citep{NNiT2026}.

\subsubsection{2.3 Model Zoo Analysis}

The Unterthiner MNIST zoo (29,997 models) provides a standard benchmark for property prediction. Peebles et al. (2022) demonstrated that diffusion Transformers applied to weight checkpoints can achieve weight-space generative modeling, confirming that Transformer architectures are viable for weight-space tasks.

\subsubsection{2.4 Positioning}

Prior work on weight-space property prediction (Unterthiner et al., 2020; Eilertsen et al., 2020) does not evaluate under permutation stress or compare against equivariant alternatives. Prior work on equivariant weight-space encoders (Zhou et al., 2023; Navon et al., 2023) does not test on model zoo property prediction. This work provides the controlled comparison that neither line has conducted.
